\subsection{正则空间}

命题 11. 拓扑空间 X 的下列性质等价：

(O \(_{III}\) ) X 的任一点的诸闭邻域组成的集合是该点的一个基本邻域系。

\((O_{\mathrm{III}}^{\prime})\) 给定闭子集 \(\mathbf{F}\)（\(\mathbf{X}\) 的）与任一点 \(x \notin \mathbf{F}\)，存在 \(x\) 的一个邻域与 \(\mathbf{F}\) 的一个邻域，二者不相交。

\((\mathrm{O}_{\mathrm{III}}) \Longrightarrow (\mathrm{O}_{\mathrm{III}}^{\prime})\)：若 \(\mathbf{F}\) 是闭的且 \(x \notin \mathbf{F}\)，则存在一个闭邻域 \(\mathbf{V}\)（即 \(x\) 的闭邻域），它包含于邻域 \(\complement \mathbf{F}\)（即 \(x\) 的邻域）中；\(\mathbf{V}\) 与 \(\complement \mathbf{V}\) 分别是 \(x\) 与 \(\mathbf{F}\) 的邻域，且没有公共点。

\((\mathrm{O}_{\mathrm{III}}^{\prime}) \Longrightarrow (\mathrm{O}_{\mathrm{III}})\)：若 \(W\) 是 \(x \in X\) 的一个开邻域，则存在邻域 \(U\)（即 \(x\) 的邻域）与邻域 \(V\)（即 \(\complement W\) 的邻域），二者不相交，因此 \(\overline{U} \subset W\)。

定义 2. 拓扑空间称为正则的，如果它是豪斯多夫的且满足公理 (O \(_{III}\) )；这时它的拓扑称为正则拓扑。

离散空间是正则的。* 我们将在 § 9 中看到，每个局部紧空间（特别地，实直线 R）都是正则的。*

命题 12. 正则空间的每个子空间都是正则的。

设 A 是正则空间 X 的一个子空间。既然 X 是豪斯多夫的，A 亦然（第 2 目）；另一方面，点 \(x \in A\) 关于 A 的每个邻域都具有形状 \(V \cap A\)，其中 V 是 x 在 X 中的邻域。既然 X 是正则的，就存在 x 在 X 中的一个邻域 W，它在 X 中是闭的且包含于 V；于是 \(W \cap A\) 是 x 在 A 中的邻域，在 A 中是闭的且包含于 \(V \cap A\)。由此即得结论。反之：

命题 13. 若拓扑空间 X 的每一点 x 都有一个闭邻域，它是 X 的正则子空间，则 X 是正则的。

由第 2 目命题 6，X 是豪斯多夫的。设 x 是 X 的任一点，V 是 x 的一个闭的正则邻域。若 U 是 x 的包含于 V 中的任一邻域，则 U 是 x 关于 V 的邻域；于是由假设，存在 x 在 V 中的一个邻域 W，它在 V 中是闭的且包含于 U。但既然 V 是 x 在 X 中的邻域，W 就是 x 在 X 中的邻域；又既然 V 在 X 中是闭的，W 就在 X 中是闭的。

注. 1) 存在这样的非豪斯多夫空间，其中每一点都有一个正则邻域（习题 7）。

\begin{enumerate}
\def\labelenumi{\arabic{enumi})}
\setcounter{enumi}{1}
\item
  存在豪斯多夫但非正则的空间（习题 20）。
\item
  比正则拓扑更细的拓扑未必是正则的（习题 20）。
\end{enumerate}

